The first property is called the center of the data.
The second is called the spread.
This chapter defines three measures of center, two measures
of spread, and examines the relationship between them.
The choice of measure is not arbitrary: different measures
reveal different features of the data, and the appropriate
choice depends on the structure of the data set itself.

%--------------------------------------------------------------------
\section{Measures of Center}
%--------------------------------------------------------------------

A \textbf{measure of center} is a single number that summarizes
the location of a data set on the number line.
Three measures are in common use: the mean, the median,
and the mode.

\subsection*{The Mean}

\begin{definition}{Mean}
The \textbf{mean} of a data set is the sum of all values
divided by the number of values.
\[
\bar{x} = \frac{\displaystyle\sum_{i=1}^{n} x_i}{n}
= \frac{x_1 + x_2 + \cdots + x_n}{n}
\]
where $n$ is the number of values and $\bar{x}$ is read
``$x$-bar.''
\end{definition}

\noindent
The mean is the most commonly used measure of center.
It uses every value in the data set.

\begin{example}
The high temperatures (in degrees Fahrenheit) recorded on
seven consecutive days were:
\[ 58,\; 63,\; 61,\; 67,\; 70,\; 65,\; 62. \]
\[
\bar{x} = \frac{58+63+61+67+70+65+62}{7}
= \frac{446}{7} \approx 63.7.
\]
The mean high temperature for the week was approximately
$63.7^\circ$F.
\end{example}

\subsection*{The Median}

\begin{definition}{Median}
The \textbf{median} of a data set is the middle value when
the values are arranged in order from smallest to largest.
When the number of values is even, the median is the mean
of the two middle values.
\end{definition}

\noindent
The median divides the data set in half: exactly half the
values fall below the median and half fall above it.

\begin{example}
Find the median of the temperature data from Example 3.1.1.

Arranging in order:
\[ 58,\; 61,\; 62,\; \mathbf{63},\; 65,\; 67,\; 70. \]
With seven values, the middle value is in position 4.
The median is $\mathbf{63}^\circ$F.
\end{example}

\begin{example}[Even number of values]
Find the median of the following eight quiz scores:
\[ 72,\; 85,\; 91,\; 78,\; 88,\; 65,\; 90,\; 80. \]

Arranging in order:
\[ 65,\; 72,\; 78,\; \mathbf{80},\; \mathbf{85},\;
   88,\; 90,\; 91. \]
The two middle values (positions 4 and 5) are $80$ and $85$.
\[
\text{Median} = \frac{80 + 85}{2} = \frac{165}{2} = 82.5.
\]
\end{example}

\subsection*{The Mode}

\begin{definition}{Mode}
The \textbf{mode} of a data set is the value that appears
most frequently.
A data set may have one mode, more than one mode, or no mode
if every value appears exactly once.
\end{definition}

\begin{example}
Find the mode of each data set.
\begin{enumerate}[label=(\alph*), itemsep=4pt]
  \item $3,\; 5,\; 5,\; 7,\; 8,\; 9,\; 5$: the value $5$
    appears three times. The mode is $5$.
  \item $2,\; 4,\; 4,\; 6,\; 6,\; 8$: both $4$ and $6$
    appear twice. This data set has two modes: $4$ and $6$.
  \item $10,\; 20,\; 30,\; 40$: every value appears once.
    There is no mode.
\end{enumerate}
\end{example}

\subsection*{Choosing a Measure of Center}

The mean uses every value in its calculation.
A single extreme value pulls the mean toward it without
affecting the median.
When a data set contains extreme values, the median is often
a more representative description of center.

\begin{example}[Effect of an extreme value]
Six employees at a small company earn the following annual
salaries (in thousands of dollars):
\[ 38,\; 42,\; 45,\; 47,\; 50,\; 310. \]

\[
\text{Mean} = \frac{38+42+45+47+50+310}{6}
= \frac{532}{6} \approx \$88{,}700.
\]

Ordered data: $38, 42, 45, 47, 50, 310$.
The two middle values are $45$ and $47$.
\[
\text{Median} = \frac{45+47}{2} = \$46{,}000.
\]

The mean of \$88{,}700 is pulled upward by the one salary
of \$310{,}000. The median of \$46{,}000 reflects what a
typical employee earns. Five of the six employees earn less
than the mean.
\end{example}

\begin{exercises}
  \item Find the mean, median, and mode of each data set.
    \begin{enumerate}[label=(\alph*), itemsep=5pt]
      \item $4,\; 7,\; 7,\; 9,\; 11,\; 13,\; 7,\; 6$
      \item $22,\; 18,\; 30,\; 22,\; 25,\; 18,\; 29,\;
             22,\; 31$
      \item $0.5,\; 1.2,\; 0.8,\; 1.5,\; 0.9,\; 1.1$
    \end{enumerate}

  \item The following data records the number of customer
    complaints received by a store over ten days:
    \[ 3,\; 0,\; 5,\; 2,\; 1,\; 4,\; 2,\; 8,\; 2,\; 3. \]
    \begin{enumerate}[label=(\alph*), itemsep=4pt]
      \item Find the mean, median, and mode.
      \item On day 8, a particularly difficult situation led
        to 8 complaints. Describe the effect of removing this
        value from the data set on each of the three measures.
    \end{enumerate}

  \item A real estate agent reports that the mean home price
    in a neighborhood is \$420{,}000, but the median is
    \$285{,}000.
    In two complete sentences, explain what this difference
    tells a potential buyer about the distribution of home
    prices in the neighborhood.

  \item Seven students took a test. Six of them scored
    $72, 78, 84, 80, 76, 88$.
    The mean score for all seven students is $80$.
    What did the seventh student score?
    Show your work.
\end{exercises}

%--------------------------------------------------------------------
\section{Measures of Spread: Range}
%--------------------------------------------------------------------

Two data sets can have the same mean but very different
distributions of values around that mean.
A measure of spread quantifies how much the values vary.

\begin{definition}{Range}
The \textbf{range} of a data set is the difference between
the maximum and minimum values.
\[
\text{Range} = \text{maximum} - \text{minimum}
\]
\end{definition}

\noindent
The range is easy to compute and gives a quick sense of the
full extent of the data.
It uses only two values and is sensitive to extreme
observations.

\begin{example}
Two groups of students each took the same exam.

Group A scores: $62, 65, 68, 70, 72, 74, 78$.
\[
\text{Mean}_A = 69.9, \qquad
\text{Range}_A = 78 - 62 = 16.
\]

Group B scores: $45, 58, 65, 72, 75, 80, 95$.
\[
\text{Mean}_B = 70.0, \qquad
\text{Range}_B = 95 - 45 = 50.
\]

The two groups have nearly identical means but very different
spreads. Group A's scores are clustered close to the center.
Group B's scores are spread widely.
\end{example}

\begin{exercises}
  \item Find the range of each data set.
    \begin{enumerate}[label=(\alph*), itemsep=4pt]
      \item $14,\; 22,\; 9,\; 31,\; 18,\; 27$
      \item $-5,\; 0,\; 3,\; -2,\; 8,\; -7,\; 1$
      \item $0.4,\; 1.1,\; 0.7,\; 1.8,\; 0.3$
    \end{enumerate}

  \item Two delivery drivers recorded the following number
    of deliveries per day over five days.
    \begin{center}
    \renewcommand{\arraystretch}{1.4}
    \begin{tabular}{lrrrrr}
    \toprule
    & Day 1 & Day 2 & Day 3 & Day 4 & Day 5 \\
    \midrule
    Driver A & 12 & 14 & 11 & 13 & 15 \\
    Driver B & 8  & 18 & 10 & 16 & 13 \\
    \bottomrule
    \end{tabular}
    \end{center}
    \begin{enumerate}[label=(\alph*), itemsep=4pt]
      \item Compute the mean and range for each driver.
      \item Which driver is more consistent?
        Justify your answer using the range.
    \end{enumerate}

  \item A data set has a range of $40$ and a minimum value
    of $15$. What is the maximum value?
\end{exercises}

%--------------------------------------------------------------------
\section{Measures of Spread: Sample Standard Deviation}
%--------------------------------------------------------------------

The range measures spread using only the two most extreme
values.
The \textbf{sample standard deviation} uses every value in the
